% Actividad Semana 4: proceso de analítica y BI con datos del SINPAD del INDECI, región Puno
% Compilar:  tectonic -X compile informe.tex
\documentclass[12pt,a4paper]{article}
\usepackage[spanish,es-noshorthands,es-tabla,es-nodecimaldot]{babel}
\usepackage{fontspec}
\setmainfont{texgyretermes}[Extension=.otf,UprightFont=*-regular,BoldFont=*-bold,ItalicFont=*-italic,BoldItalicFont=*-bolditalic]
\setsansfont{texgyreheros}[Extension=.otf,UprightFont=*-regular,BoldFont=*-bold,ItalicFont=*-italic,BoldItalicFont=*-bolditalic]
\setmonofont{texgyrecursor}[Extension=.otf,UprightFont=*-regular,BoldFont=*-bold,ItalicFont=*-italic,BoldItalicFont=*-bolditalic,Scale=0.95]
\usepackage[a4paper,margin=2.5cm]{geometry}
\usepackage{setspace}
\usepackage{graphicx}
\usepackage{booktabs,tabularx,array,multirow}
\usepackage[table]{xcolor}
\usepackage{caption}
\usepackage{float}
\usepackage{enumitem}
\usepackage{amsmath}
\usepackage[round,authoryear]{natbib}
\usepackage{titlesec}
\usepackage{fancyhdr}
\usepackage[most]{tcolorbox}
\usepackage{tikz}
\usetikzlibrary{arrows.meta,positioning}
\usepackage[hidelinks]{hyperref}
\usepackage{xurl}

\graphicspath{{../figuras/}{../capturas/informe/}}
\hypersetup{colorlinks=true,linkcolor=black,citecolor=black,urlcolor=blue}

\definecolor{marino}{RGB}{31,36,112}
\definecolor{azul}{RGB}{42,120,214}
\definecolor{fondo}{RGB}{243,247,254}
\definecolor{naranja}{RGB}{235,104,52}
\definecolor{fondonar}{RGB}{255,246,240}

\onehalfspacing
\setlength{\emergencystretch}{3em}
\setlength{\parindent}{1.27cm}
\setlength{\parskip}{2pt}
\captionsetup{font=small,labelfont=bf,justification=justified,singlelinecheck=true}
\captionsetup[figure]{name=Figura}
\titleformat{\section}{\large\bfseries}{\thesection.}{0.6em}{}
\titleformat{\subsection}{\normalsize\bfseries}{\thesubsection}{0.6em}{}
\titlespacing*{\section}{0pt}{14pt}{6pt}
\titlespacing*{\subsection}{0pt}{10pt}{4pt}
\setlist{itemsep=2pt,topsep=4pt}
\newcolumntype{L}[1]{>{\raggedright\arraybackslash}p{#1}}
\newcolumntype{Y}{>{\raggedright\arraybackslash}X}

\setlength{\headheight}{15pt}
\pagestyle{fancy}
\fancyhf{}
\fancyhead[L]{\small Actividad Semana 4: proceso de analítica y BI}
\fancyhead[R]{\small Emergencias en la región Puno}
\fancyfoot[C]{\small\thepage}
\renewcommand{\headrulewidth}{0.4pt}

\newcommand{\fig}[3]{% 1 archivo, 2 ancho, 3 pie y etiqueta
  \begin{figure}[H]\centering\includegraphics[width=#2\textwidth]{#1}#3\end{figure}}

% Caja de interpretación en lenguaje sencillo
\newtcolorbox{lectura}{enhanced,breakable,colback=fondo,colframe=azul,boxrule=0pt,leftrule=3pt,arc=2pt,
  left=8pt,right=8pt,top=5pt,bottom=5pt,before skip=6pt,after skip=8pt,fontupper=\small\setlength{\parindent}{0pt},
  title={Interpretación del resultado},coltitle=azul,colbacktitle=fondo,fonttitle=\small\bfseries,
  attach title to upper={\par\smallskip}}
\newtcolorbox{sintesis}[1]{enhanced,breakable,colback=white,colframe=marino,boxrule=0.8pt,arc=3pt,
  left=10pt,right=10pt,top=6pt,bottom=6pt,colbacktitle=marino,coltitle=white,fonttitle=\bfseries,title={#1}}
\newtcolorbox{alerta}{enhanced,breakable,colback=fondonar,colframe=naranja,boxrule=0pt,leftrule=3pt,arc=2pt,
  left=8pt,right=8pt,top=5pt,bottom=5pt,before skip=6pt,after skip=8pt,fontupper=\small\setlength{\parindent}{0pt}}

\begin{document}

% =============================================================================
\thispagestyle{plain}
\begin{center}
{\large\bfseries Emergencias por peligros naturales en la región Puno entre 2003 y 2026:\\
analítica descriptiva, diagnóstica, predictiva y prescriptiva con datos del SINPAD\par}
\vspace{10pt}
{Mary Luz Nina Palacios\par}
{\small Universidad Nacional del Altiplano, Escuela de Posgrado. Maestría en Informática, curso Business
Intelligence and Data Analytics. Docente: Dr. Fred Torres Cruz\par}
{\small Actividad Semana 4: Analítica, modelos y optimización. Puno, noviembre de 2026\par}
\end{center}
\vspace{4pt}

% =============================================================================
\section*{Resumen}
\noindent Este informe documenta el proceso completo de analítica e inteligencia de negocios aplicado a las 6\,593
emergencias que el Sistema de Información Nacional para la Respuesta y Rehabilitación del Instituto Nacional de
Defensa Civil registró en la región Puno entre enero de 2003 y junio de 2026. La revisión de calidad detectó
subregistro en 2014 y 2019, la pérdida progresiva del registro de la ayuda humanitaria y un error de origen en la
columna de viviendas destruidas del archivo de 2026. El análisis exploratorio muestra que las heladas representan
el 29\,\% de las emergencias, pero concentran el 77\,\% de las personas impactadas y el 76\,\% de los camélidos
perdidos, y que el riesgo se reparte en dos temporadas: lluvias, de diciembre a marzo, y heladas, de mayo a agosto.
Se planteó predecir, con la información disponible al momento del reporte, si una emergencia afectará a 100 o más
personas. De ocho modelos comparados con validación cruzada temporal, Random Forest obtuvo el mejor desempeño,
con un ROC AUC de 0,882. Sus hiperparámetros se optimizaron con Optuna, scikit-optimize, Hyperopt y búsqueda
aleatoria bajo el mismo presupuesto; todas las estrategias alcanzaron valores entre 0,890 y 0,891, y Optuna fue la
más rápida gracias a la poda. En el periodo de prueba, de 2022 a 2026, el modelo optimizado alcanzó un ROC AUC de
0,821, frente a 0,800 con los valores por defecto. Con el modelo se construyó una matriz de riesgo que anticipa
unas 76 emergencias graves por año, encabezadas por Melgar, Azángaro y Puno. La analítica prescriptiva propone
priorizar las evaluaciones de daños según el riesgo estimado, lo que permite atender el 54\,\% de los casos graves
evaluando solo el 20\,\% de los reportes, y un modelo de programación lineal que distribuye los kits de ayuda entre
provincias antes de cada temporada. Los resultados se integran en un tablero interactivo publicado en línea y el
proceso completo es reproducible desde un repositorio público.

\medskip\noindent\textbf{Palabras clave:} inteligencia de negocios, gestión del riesgo de desastres, SINPAD,
Random Forest, optimización de hiperparámetros, Optuna, programación lineal, tablero interactivo.

\begin{sintesis}{Síntesis para la toma de decisiones}
\small
\begin{itemize}[leftmargin=14pt]
\item \textbf{El problema central de Puno son las heladas.} Ocurren menos que otros peligros, pero dejan casi
ocho de cada diez personas impactadas y tres de cada cuatro alpacas y llamas perdidas.
\item \textbf{El riesgo tiene calendario.} Las lluvias golpean entre diciembre y marzo, y las heladas entre mayo y
agosto. Abril y noviembre son las ventanas para prepararse.
\item \textbf{La gravedad se puede anticipar.} Con solo saber qué peligro ocurrió, dónde y cuándo, y cómo le fue antes a
ese lugar, un modelo distingue los casos graves de los leves en ocho de cada diez comparaciones.
\item \textbf{Se esperan unas 76 emergencias graves por año.} Melgar, Azángaro y Puno concentran el mayor riesgo.
\item \textbf{Dos decisiones concretas.} Evaluar primero los reportes de mayor riesgo, lo que duplica la cobertura
de los casos graves con el mismo personal, y ubicar al menos el 75\,\% de los kits de ayuda en las provincias antes
de cada temporada.
\end{itemize}
\textbf{Tablero interactivo:} \url{https://mlnp.shinyapps.io/emergencias-puno/}\\
\textbf{Repositorio con datos, código y resultados:} \url{https://github.com/mlnp-dev/emergencias-puno-indeci}
\end{sintesis}

% =============================================================================
\section{Introducción}

Puno es una de las regiones del Perú que más emergencias registra cada año. En el altiplano, por encima de los
3\,800 metros, las heladas y las nevadas de junio y julio enferman a la población, causan la muerte de alpacas y
llamas y deterioran los pastos; entre diciembre y marzo, las lluvias intensas desbordan ríos, inundan viviendas y
activan huaicos en la selva alta de Sandia y Carabaya. Cada uno de estos eventos queda registrado en el Sistema de
Información Nacional para la Respuesta y Rehabilitación, conocido como SINPAD, que administra el Instituto Nacional
de Defensa Civil, junto con la Evaluación de Daños y Análisis de Necesidades que elaboran los gobiernos locales.

Ese registro constituye una fuente valiosa para la inteligencia de negocios en el sector público. No basta con
contar cuántas emergencias ocurrieron: interesa saber dónde y cuándo se concentran, qué hace que una emergencia
termine afectando a cientos de personas y, sobre todo, cómo usar esa información para decidir antes de la
siguiente temporada. Esta progresión corresponde a los cuatro niveles de la analítica de negocios, descriptivo,
diagnóstico, predictivo y prescriptivo, que proponen \citet{delen2018} y que \citet{lepenioti2020} desarrollan
para el nivel prescriptivo.

La pregunta de negocio que guía el trabajo es la siguiente: \textbf{¿se puede estimar, en el momento en que un
distrito reporta una emergencia, la probabilidad de que llegue a afectar a 100 o más personas, y cómo debe usarse
esa estimación para asignar los equipos de evaluación y la ayuda humanitaria?}

\noindent\textbf{Objetivo general.} Desarrollar un proceso integral de analítica e inteligencia de negocios sobre
las emergencias registradas en la región Puno que transforme el registro histórico del SINPAD en información útil
para la gestión del riesgo de desastres.

\noindent\textbf{Objetivos específicos.}
\begin{enumerate}[leftmargin=22pt]
\item Seleccionar, depurar y documentar un conjunto de datos local y abierto, evaluando la calidad de su registro.
\item Describir el comportamiento de las emergencias según el peligro, el territorio y el calendario.
\item Comparar modelos de clasificación que anticipen el alto impacto de una emergencia con validación temporal.
\item Optimizar los hiperparámetros del mejor modelo con las herramientas revisadas en clase y comparar su eficiencia.
\item Comunicar los resultados mediante las analíticas descriptiva, diagnóstica, predictiva y prescriptiva, y
mediante un tablero interactivo de acceso público.
\end{enumerate}

El informe se ordena en función de esos objetivos. La sección 2 resume el marco de referencia y la sección 3
describe la metodología. Las secciones 4 a 8 desarrollan los cinco pasos del encargo: selección del conjunto de
datos, análisis exploratorio, selección de modelos, optimización y \emph{reporting}. La sección 9 presenta el
tablero interactivo, y las secciones finales discuten los resultados, sus limitaciones y las recomendaciones.

% =============================================================================
\section{Marco de referencia}

La analítica de negocios se entiende como una progresión de preguntas. La analítica descriptiva responde qué
pasó; la diagnóstica, por qué pasó; la predictiva, qué puede pasar, y la prescriptiva, qué conviene hacer
\citep{delen2018}. Este último nivel es el que más valor aporta a la decisión, porque combina las predicciones con
modelos de optimización o reglas de decisión que recomiendan acciones concretas \citep{lepenioti2020}.

En la gestión de desastres, el aprendizaje automático se ha aplicado a la alerta temprana, la estimación de daños
y la asignación de recursos. Las revisiones de \citet{sun2020}, \citet{linardos2022} y \citet{kyrkou2023}
coinciden en que los modelos de clasificación basados en árboles ofrecen un buen equilibrio entre precisión e
interpretabilidad, y en que la principal dificultad no es el algoritmo sino la calidad y la disponibilidad de los
registros históricos. \citet{ghaffarian2023} añaden que, para que estos modelos se usen en la gestión pública, sus
resultados deben poder explicarse a quienes toman las decisiones, lo que justifica el énfasis de este trabajo en
la interpretación de cada resultado.

La optimización de hiperparámetros busca la configuración de un modelo que maximiza su desempeño fuera de la
muestra. \citet{bischl2023} distinguen entre la búsqueda aleatoria, que sirve como referencia, y los métodos
basados en modelos sustitutos, como la optimización bayesiana con procesos gaussianos \citep{shahriari2016} y el
estimador de Parzen estructurado en árbol \citep{bergstra2015}, que aprenden de las evaluaciones anteriores para
proponer las siguientes. Optuna \citep{akiba2019} incorpora además la poda, que abandona las configuraciones poco
prometedoras antes de terminar de evaluarlas.

Por último, ubicar la ayuda humanitaria antes de que ocurra el desastre es un problema clásico de la logística
humanitaria. Los modelos de pre posicionamiento deciden cuánto almacenar y dónde, equilibrando la cobertura de la
demanda incierta con los costos de transporte \citep{balcik2008,rawls2010}. La revisión de \citet{sabbaghtorkan2020}
muestra que la programación lineal es la herramienta más utilizada para este fin y que una de las brechas
pendientes es alimentar estos modelos con estimaciones de demanda basadas en datos, que es precisamente lo que
aporta la parte predictiva de este trabajo.

% =============================================================================
\section{Metodología}

El trabajo siguió un proceso de inteligencia de negocios de seis etapas, desde la fuente de datos hasta el
tablero de decisión, como resume la Figura~\ref{fig:proceso}. Cada etapa produce un insumo para la siguiente y todas
quedaron programadas, de modo que cualquier cifra del informe puede reproducirse a partir de los archivos
originales.

\begin{figure}[H]\centering
\begin{tikzpicture}[node distance=4mm and 2.6mm,
  caja/.style={draw=marino,fill=fondo,rounded corners=3pt,align=center,font=\scriptsize,minimum height=16mm,text width=21.5mm,line width=0.6pt},
  flecha/.style={-{Stealth[length=2mm]},color=marino,line width=0.8pt}]
\node[caja] (a) {\textbf{1. Fuente}\\SINPAD del INDECI, datos abiertos};
\node[caja,right=of a] (b) {\textbf{2. Preparación}\\Limpieza, calidad y variables};
\node[caja,right=of b] (c) {\textbf{3. Exploración}\\Descriptiva y diagnóstica};
\node[caja,right=of c] (d) {\textbf{4. Modelado}\\Ocho modelos, validación temporal};
\node[caja,right=of d] (e) {\mbox{\textbf{5. Optimización}}\\Optuna, skopt, Hyperopt y PuLP};
\node[caja,right=of e,fill=marino,text=white] (f) {\textbf{6. Reporting}\\Informe y tablero en Shiny};
\foreach \x/\y in {a/b,b/c,c/d,d/e,e/f} \draw[flecha] (\x) -- (\y);
\end{tikzpicture}
\caption{Proceso de analítica e inteligencia de negocios seguido en el trabajo.}\label{fig:proceso}
\end{figure}

La Tabla~\ref{tab:herramientas} detalla las herramientas utilizadas en cada etapa. El análisis se programó en
Python dentro de un cuaderno Jupyter y el tablero se desarrolló en R con Shiny, de forma que la parte analítica y la
parte de comunicación quedan separadas pero comparten los mismos resultados exportados.

\begin{table}[H]\centering\small
\caption{Herramientas utilizadas en cada etapa del proceso}\label{tab:herramientas}
\begin{tabularx}{\textwidth}{L{3.1cm}Y}
\toprule
Etapa & Herramientas \\
\midrule
Preparación y exploración & Python 3, pandas, NumPy, SciPy, Matplotlib y Seaborn \\
Modelado & scikit-learn \citep{pedregosa2011}, XGBoost \citep{chen2016} y LightGBM \citep{ke2017} \\
Optimización de hiperparámetros & Optuna \citep{akiba2019}, scikit-optimize y Hyperopt \citep{bergstra2015} \\
Optimización de decisiones & PuLP con el resolvedor CBC para programación lineal \\
Reporting & \LaTeX\ para el informe; R, Shiny, plotly y leaflet para el tablero, publicado en shinyapps.io \\
Reproducibilidad & Jupyter Notebook y repositorio público en GitHub \\
\bottomrule
\end{tabularx}
\end{table}

% =============================================================================
\section{Paso 1: selección del conjunto de datos}\label{sec:datos}

\subsection{Fuente y justificación}

Se utilizó el conjunto \emph{Emergencias y daños a nivel nacional por departamento}, publicado por el INDECI en la
Plataforma Nacional de Datos Abiertos con licencia ODC BY \citep{indeci2026}. Está formado por dos tablas, la serie
histórica de 2003 a 2025, con 142\,139 registros, y el primer semestre de 2026, con 9\,352 registros, acompañadas
de un diccionario de datos y un documento de metadatos. Cada fila corresponde a una emergencia ocurrida en un
distrito y contiene 49 columnas: código SINPAD, fecha, ubicación, peligro y tipo de peligro, daños en personas,
viviendas, locales educativos y de salud, vías, canales de riego, cultivos y ganado, así como el peso y el costo de
la ayuda humanitaria entregada.

Puno ocupa el séptimo lugar nacional por número de emergencias registradas, con el 4,4\,\% del total, después de
Apurímac, Huancavelica, Cusco, Ayacucho, Cajamarca y Huánuco. Sus 6\,593 registros, distribuidos en 13 provincias y
109 distritos, ofrecen volumen suficiente para modelar. El tema cumple, además, la condición de ser un conjunto de
datos único dentro del aula, ya que las otras fuentes anunciadas corresponden a la producción de quinua, los
visitantes a sitios turísticos, la producción de trucha y las ventas de una distribuidora de bebidas.

\subsection{Preparación}

La preparación resolvió cinco problemas detectados al leer los archivos:

\begin{enumerate}
\item \textbf{Nombres de columnas distintos.} El archivo de 2026 renombra varias columnas, aunque conserva su orden.
Se asignaron nombres cortos comunes a las 49 columnas.
\item \textbf{Dos formatos de fecha.} La serie histórica usa día, mes y año con cuatro dígitos, mientras que la de 2026
usa mes, día y año con dos dígitos. Cada fecha se interpretó según su longitud y se contrastó con las columnas de
año y mes, sin encontrar fechas inválidas ni discrepancias.
\item \textbf{Tildes y catálogo de peligros cambiante.} En 2026 aparecen Azángaro, San Román y Huancané con tilde, que
sin corregir se contarían como provincias nuevas, y el catálogo de peligros cambia de nombres entre años. Se
normalizaron las tildes y los 34 nombres de peligro se agruparon en diez familias: heladas y bajas temperaturas,
que incluye las nevadas; lluvias, granizo y tormentas; inundaciones; vientos fuertes; incendios urbanos; incendios
forestales; sequía; movimientos en masa; sismos y vulcanismo; y otros.
\item \textbf{Error en las viviendas destruidas de 2026.} En el 78\,\% de las filas de Puno de 2026 con viviendas
destruidas, el valor coincide con el número de personas afectadas, y el mismo patrón aparece en el 57\,\% de las
filas a nivel nacional. Ello arrojaba 44\,930 viviendas destruidas en un semestre, frente a 17\,382 en los 23 años
anteriores. Al tratarse de un error de origen, esa columna se anuló para 2026.
\item \textbf{Variables derivadas.} Se definieron las \emph{personas impactadas} como la suma de fallecidos,
desaparecidos, lesionados, damnificados y afectados, y la \emph{emergencia de alto impacto} como aquella con 100 o
más personas impactadas. Con ese umbral, 1\,876 emergencias, el 28,5\,\% del total, son de alto impacto.
\end{enumerate}

\begin{lectura}
Antes de analizar fue necesario poner en orden los datos, porque el INDECI cambió con los años la forma de anotar
fechas, nombres y peligros. Sin esa corrección, una misma provincia habría aparecido dos veces y el primer semestre
de 2026 habría mostrado más viviendas destruidas que los 23 años anteriores juntos, lo que es imposible. La base
final tiene 6\,593 emergencias, y casi tres de cada diez afectaron a 100 personas o más.
\end{lectura}

% =============================================================================
\section{Paso 2: análisis exploratorio de datos}\label{sec:eda}

\subsection{Calidad del registro}

La Figura~\ref{fig:anio} muestra dos problemas que condicionan todo el análisis. El primero es el subregistro de
algunos años: 2019 tiene solo 43 emergencias y 2014 solo 83, frente a un promedio de 248 por año entre 2003 y 2021.
No fueron años sin heladas ni lluvias, sino años en que el sistema registró poco. El segundo es que el registro de
la ayuda humanitaria se fue perdiendo. Entre 2003 y 2011 más del 90\,\% de las emergencias tenía la ayuda anotada;
desde 2015 la proporción baja del 40\,\% y en 2019 y entre 2024 y 2026 es nula. Por ello, la ayuda entregada no se
usó como variable objetivo. Además, entre 2014 y 2019 creció la proporción de emergencias registradas sin ninguna
persona impactada, hasta llegar al 57\,\% en 2017. Desde 2022 el registro vuelve a crecer, con un promedio de 403
emergencias por año entre 2022 y 2025.

\fig{fig01_eventos_por_anio}{0.95}{\caption{Emergencias registradas por año en Puno y proporción con ayuda
humanitaria registrada. El año 2026 corresponde solo al primer semestre.}\label{fig:anio}}

\begin{alerta}
\textbf{Advertencia sobre los datos.} Un año con pocas emergencias en la base no fue necesariamente un año tranquilo.
En 2014 y 2019 se registró mucho menos de lo que ocurrió, y desde 2015 casi no se anota la ayuda entregada. Por eso
estas cifras no deben leerse como si la ayuda hubiera desaparecido, sino como un problema del llenado del sistema.
\end{alerta}

\subsection{Distribución del impacto}

El impacto es muy asimétrico, como se aprecia en la Figura~\ref{fig:dist}. El 22,5\,\% de las emergencias no tiene
ninguna persona impactada, la mediana es de 9 personas y el percentil 90 supera las 1\,100. Unas pocas emergencias
alcanzan decenas de miles de personas; las dos mayores de la serie ocurrieron en Ilave, provincia de El Collao, en
enero de 2016, una por sequía, con 54\,590 personas, y otra por bajas temperaturas, con 50\,305. La suma de todas
las emergencias asciende a 2,76 millones de personas impactadas, pero esa cifra debe leerse como personas por
evento: una familia afectada por las heladas en varias temporadas se cuenta varias veces, por lo que no es
comparable con la población de la región, de alrededor de 1,2 millones.

\fig{fig02_distribucion_personas}{0.72}{\caption{Distribución de las personas impactadas por emergencia en escala
logarítmica y umbral de alto impacto.}\label{fig:dist}}

\begin{lectura}
La mayoría de emergencias son pequeñas: la mitad afecta a 9 personas o menos. Sin embargo, unas pocas afectan a miles
y son las que concentran casi todo el daño. Por eso conviene identificar a tiempo cuáles serán graves, que es el
objetivo del modelo desarrollado en el paso 3.
\end{lectura}

\subsection{Frecuencia frente a impacto según el peligro}

Lo más frecuente no es lo que más daño causa, como muestra la Figura~\ref{fig:peligro}. Las heladas y bajas
temperaturas representan el 29\,\% de las emergencias, pero reúnen el 77\,\% de las personas impactadas y el 76\,\%
de los camélidos perdidos, y el 66\,\% de ellas supera las 100 personas. Los incendios urbanos son el segundo
peligro más frecuente, con 1\,170 casos, pero casi nunca alcanzan el alto impacto, apenas en el 2\,\% de los casos.
Las inundaciones y la sequía ocurren con menor frecuencia, aunque un tercio de sus emergencias son de alto impacto;
la sequía casi no deja damnificados, pero explica el 20\,\% de los camélidos perdidos y buena parte de los cultivos
dañados.

\fig{fig03_frecuencia_vs_impacto}{1}{\caption{Número de emergencias y personas impactadas por familia de
peligro en Puno entre 2003 y 2026.}\label{fig:peligro}}

\begin{lectura}
Si solo se contara cuántas emergencias hay, se pensaría que los incendios urbanos son un problema tan grande como las
heladas. Al medir a cuántas personas afectan, la conclusión cambia por completo: las heladas causan casi ocho de cada
diez personas impactadas. La prioridad de la gestión del riesgo en Puno debe ser, por tanto, la temporada de frío.
\end{lectura}

\subsection{Estacionalidad}

La región presenta dos temporadas de emergencias muy marcadas, que se observan en la Figura~\ref{fig:mes}. Entre
diciembre y marzo ocurre el 78\,\% de las emergencias por lluvias, granizo y tormentas y el 89\,\% de las
inundaciones; febrero es el mes con más emergencias del año. Entre mayo y agosto ocurre el 64\,\% de las
emergencias por bajas temperaturas, y solo junio y julio reúnen el 46\,\%. Al final de la época seca, de agosto a
octubre, se concentra el 71\,\% de los incendios forestales, junto con los vientos fuertes. Abril, mayo y diciembre
son los meses más tranquilos y marcan las ventanas para preparar la respuesta.

\fig{fig05_estacionalidad}{1}{\caption{Emergencias por mes y familia de peligro.}\label{fig:mes}}

\subsection{Distribución territorial}

Las heladas dominan en casi todas las provincias altiplánicas, como se observa en la Figura~\ref{fig:prov}. Sandia,
la provincia de selva alta, presenta otro perfil, ya que el 56\,\% de sus emergencias corresponde a lluvias intensas
y movimientos en masa. Moho y Chucuito tienen una proporción alta de incendios urbanos, mientras que en El Collao y
Yunguyo destacan los vientos fuertes. Por personas impactadas acumuladas sobresalen Azángaro, Puno y El Collao.
Ilave es el distrito más afectado, con unas 184 mil personas por evento en 241 emergencias, y los 15 distritos con
más personas impactadas reúnen el 34\,\% del total, según la Figura~\ref{fig:dist15}.

\fig{fig06_provincia_peligro}{1}{\caption{Perfil de peligros por provincia. El número indica las emergencias y el
color, su participación dentro de la provincia.}\label{fig:prov}}

\fig{fig07_top_distritos}{0.85}{\caption{Los 15 distritos con más personas impactadas, acumulado de 2003 a 2026.}\label{fig:dist15}}

\begin{lectura}
El riesgo no es igual en toda la región. En el altiplano el enemigo es el frío, mientras que en la selva alta de
Sandia lo son las lluvias y los deslizamientos. Además, el daño se concentra en pocos lugares: quince distritos
acumulan un tercio de todas las personas impactadas. Esto permite focalizar la prevención en lugar de repartir los
recursos por igual.
\end{lectura}

\subsection{Tendencias y daños productivos}

Las pérdidas de camélidos se concentran en pocos años: 2010, con unos 300 mil animales, y 2015, con 112 mil,
registradas como bajas temperaturas, y 2023, con 100 mil, registradas como sequía, como muestra la
Figura~\ref{fig:tend}. El año 2023 es también el de mayor superficie de cultivo dañada, con unas 272 mil hectáreas,
el 93\,\% por sequía. Las correlaciones de Spearman muestran que damnificados y afectados miden aspectos distintos:
los damnificados se asocian con las viviendas destruidas, con $\rho=0{,}79$, y los afectados con el total de
personas, con $\rho=0{,}84$. La ayuda entregada apenas se relaciona con los daños, con $\rho=0{,}36$ respecto de
las personas impactadas, lo que constituye otra señal de su registro irregular.

\fig{fig08_tendencias_y_danos}{1}{\caption{Emergencias por año de los cuatro peligros más frecuentes, a la izquierda,
y pérdidas de camélidos y cultivos, a la derecha.}\label{fig:tend}}

% =============================================================================
\section{Paso 3: selección de modelos}\label{sec:modelos}

\subsection{Planteamiento}

Cuando un distrito reporta una emergencia, todavía no se conoce cuántas personas resultarán afectadas; ese dato se
obtiene días después, con la evaluación de daños. El problema se formuló, por tanto, como una clasificación binaria:
estimar, en el primer momento, la probabilidad de que la emergencia sea de alto impacto, es decir, que afecte a 100
o más personas, para decidir a dónde enviar primero a los evaluadores y los bienes de ayuda.

Para que el modelo resulte útil en la práctica, solo se usaron variables conocidas al momento del reporte. Los daños
quedaron fuera porque son justamente lo que se quiere anticipar, y usarlos constituiría una fuga de información
\citep{kaufman2012}. Las variables se presentan en la Tabla~\ref{tab:variables}. Las de historia del lugar se
calcularon únicamente con emergencias de fechas anteriores, y las tasas se suavizaron hacia la tasa general para que
un distrito con una sola emergencia previa no reciba un valor extremo de 0 o 100\,\%.

\begin{table}[H]\centering\small
\caption{Variables predictoras, todas disponibles al momento del reporte}\label{tab:variables}
\begin{tabularx}{\textwidth}{L{3.2cm}Y}
\toprule
Grupo & Variables \\
\midrule
Qué pasó & Familia de peligro; tipo de peligro, natural o inducido por la acción humana \\
Dónde & Provincia; región natural, sierra o selva alta \\
Cuándo & Mes, codificado con seno y coseno para respetar su carácter cíclico; año \\
Historia del distrito & Emergencias previas; tasa previa de alto impacto; promedio previo de personas
impactadas en logaritmo; días desde la última emergencia \\
Historia de la provincia & Emergencias previas y tasa previa de alto impacto del mismo peligro en la provincia;
emergencias en la provincia en los 30 días anteriores \\
\bottomrule
\end{tabularx}
\end{table}

La partición respetó el orden temporal. El entrenamiento usó las emergencias de 2003 a 2021, que suman 4\,715 con un
31,2\,\% de alto impacto, y la prueba las de 2022 a 2026, que suman 1\,878 con un 21,6\,\% de alto impacto y que el
modelo no vio en ningún momento. Dentro del entrenamiento, la validación cruzada usó cinco bloques consecutivos:
cada bloque se entrena con el pasado y se valida con el periodo siguiente. Mezclar al azar emergencias de distintos
años produciría una estimación demasiado optimista \citep{bergmeir2012,cerqueira2020}.

\subsection{Modelos comparados y métricas}

Se compararon ocho clasificadores: una línea base que predice al azar con la proporción de clases, regresión
logística, árbol de decisión, $k$ vecinos más cercanos, Random Forest \citep{breiman2001}, HistGradientBoosting,
XGBoost y LightGBM. Las variables categóricas se transformaron en indicadores binarios y, para la regresión
logística y los $k$ vecinos, las numéricas se estandarizaron. En los modelos que lo permiten se compensó el
desbalance de clases. La métrica principal fue el ROC AUC, que mide la capacidad del modelo para ordenar las
emergencias según su riesgo sin depender de un umbral. Se reportan también el PR AUC, más informativo cuando las
clases están desbalanceadas \citep{saito2015}, la F1, la exactitud balanceada y el puntaje de Brier, que mide la
calibración de las probabilidades. El significado de cada métrica se explica en el glosario del Anexo A.

\subsection{Resultados}

La Tabla~\ref{tab:seleccion} y la Figura~\ref{fig:seleccion} resumen la comparación. Todos los modelos superan con
amplitud a la línea base, lo que confirma que la información disponible al momento del reporte permite anticipar el
impacto. Random Forest obtuvo el mejor ROC AUC, 0,882, la mejor F1 y la mejor calibración. Las diferencias entre
los seis primeros modelos son pequeñas frente a la variación entre bloques, cuya desviación ronda 0,04, de modo
que el periodo de validación influye más que el algoritmo. El árbol de decisión simple fue claramente inferior y el
más inestable. Random Forest pasó al paso de optimización.

\begin{table}[H]\centering\small
\caption{Selección de modelos: promedio de los cinco bloques de validación temporal, 2003 a 2021}\label{tab:seleccion}
\begin{tabular}{lcccccc}
\toprule
Modelo & ROC AUC & Desv. & PR AUC & F1 & Exact. bal. & Brier \\
\midrule
\textbf{Random Forest} & \textbf{0,882} & 0,043 & 0,772 & \textbf{0,742} & \textbf{0,809} & \textbf{0,122} \\
$k$ vecinos, $k=25$ & 0,879 & 0,036 & 0,760 & 0,714 & 0,791 & 0,125 \\
Regresión logística & 0,875 & 0,041 & \textbf{0,778} & 0,679 & 0,780 & 0,162 \\
HistGradientBoosting & 0,868 & 0,041 & 0,758 & 0,696 & 0,777 & 0,128 \\
XGBoost & 0,867 & 0,037 & 0,759 & 0,711 & 0,785 & 0,126 \\
LightGBM & 0,864 & 0,040 & 0,760 & 0,711 & 0,783 & 0,128 \\
Árbol de decisión & 0,771 & 0,114 & 0,635 & 0,625 & 0,735 & 0,173 \\
Línea base al azar & 0,507 & 0,016 & 0,316 & 0,312 & 0,507 & 0,427 \\
\bottomrule
\end{tabular}
\end{table}

\fig{fig10_seleccion_modelos}{0.82}{\caption{ROC AUC de cada modelo en los cinco bloques de validación temporal.}\label{fig:seleccion}}

\begin{lectura}
El ROC AUC puede leerse así: si se toma al azar una emergencia grave y una leve, es la probabilidad de que el modelo
asigne más riesgo a la grave. Un valor de 0,5 equivale a lanzar una moneda y 1 sería un acierto perfecto. Random
Forest acierta en 88 de cada 100 comparaciones. Que varios modelos queden muy cerca indica que el resultado es
sólido: la información del lugar, el mes y la historia del distrito sí anticipa la gravedad, sin depender de un
algoritmo en particular.
\end{lectura}

% =============================================================================
\section{Paso 4: optimización de hiperparámetros}\label{sec:opt}

\subsection{Estrategias comparadas}

De las cinco herramientas recomendadas en clase se ejecutaron tres, y se añadió la búsqueda aleatoria como punto de
comparación, según la Tabla~\ref{tab:estrategias}. Las cuatro estrategias tuvieron el mismo presupuesto de 40
evaluaciones, la misma métrica, que fue el ROC AUC promedio en los cinco bloques temporales, y el mismo espacio de
búsqueda para Random Forest: número de árboles entre 100 y 300, profundidad máxima entre 3 y 30, mínimo de casos
por hoja entre 1 y 50 en escala logarítmica y proporción de variables consideradas en cada división entre 0,1 y 1,0.

\begin{table}[H]\centering\small
\caption{Estrategias de optimización comparadas}\label{tab:estrategias}
\begin{tabularx}{\textwidth}{L{3.3cm}L{3.6cm}Y}
\toprule
Estrategia & Librería & Cómo elige la siguiente configuración \\
\midrule
Búsqueda aleatoria & Optuna, muestreador aleatorio & Al azar; sirve de referencia \\
TPE con poda & Optuna \citep{akiba2019} & Estima qué zonas del espacio dan buenos resultados y abandona a mitad de
camino las pruebas que van peor que la mediana en los primeros bloques \\
Proceso gaussiano & scikit-optimize & Optimización bayesiana clásica: ajusta un proceso gaussiano a los resultados
y elige el punto con mayor mejora esperada \citep{shahriari2016} \\
TPE clásico & Hyperopt \citep{bergstra2015} & Algoritmo TPE original, sin poda \\
\bottomrule
\end{tabularx}
\end{table}

Ray Tune y BoTorch con Ax se revisaron, pero no se ejecutaron. Ray Tune está orientado a distribuir pruebas entre
muchos procesadores o máquinas y, internamente, puede usar los mismos algoritmos de Optuna o Hyperopt; BoTorch
requiere PyTorch y está pensado para espacios continuos con evaluaciones muy costosas \citep{bischl2023}. Con
4\,715 filas y entrenamientos de pocos segundos no habrían modificado la comparación, y su instalación excedía la
memoria y el espacio disponibles en el equipo.

\subsection{Resultados}

Las cuatro estrategias mejoraron el modelo por defecto, que pasó de 0,882 a alrededor de 0,891 de ROC AUC en
validación, como muestran la Tabla~\ref{tab:opt} y la Figura~\ref{fig:conv}. Optuna obtuvo el mejor valor, en el
menor tiempo y abandonando 5 pruebas gracias a la poda. scikit-optimize llegó casi al mismo valor en solo 15
evaluaciones, lo que confirma que el proceso gaussiano aprende con rapidez cuando el espacio tiene pocas
dimensiones. Hyperopt y la búsqueda aleatoria terminaron muy cerca. La diferencia entre estrategias, menor que
0,001, es mucho más pequeña que la variación entre bloques, de manera que la elección del optimizador influyó poco
en el resultado final y bastante en el costo: Optuna necesitó un 30\,\% menos de tiempo que la búsqueda aleatoria.

\begin{table}[H]\centering\small
\caption{Comparación de estrategias de optimización con 40 evaluaciones cada una}\label{tab:opt}
\begin{tabular}{lcccc}
\toprule
Estrategia & Mejor ROC AUC & Evaluación en que lo alcanzó & Pruebas podadas & Tiempo en s \\
\midrule
Optuna, TPE con poda & \textbf{0,8911} & 31 & 5 & \textbf{236} \\
scikit-optimize, proceso gaussiano & 0,8907 & \textbf{15} & 0 & 296 \\
Hyperopt, TPE & 0,8906 & 23 & 0 & 366 \\
Búsqueda aleatoria & 0,8904 & 30 & 0 & 338 \\
\midrule
Random Forest por defecto & 0,8816 & & & \\
\bottomrule
\end{tabular}
\end{table}

\fig{fig11_convergencia_optimizadores}{1}{\caption{Mejor ROC AUC encontrado por cada estrategia según el número de
evaluaciones, a la izquierda, y según el tiempo, a la derecha.}\label{fig:conv}}

Los mejores hiperparámetros, encontrados por Optuna, fueron 158 árboles, profundidad máxima de 14, al menos 46
casos por hoja y el 71\,\% de las variables en cada división. Las cuatro estrategias coincidieron en lo esencial:
hojas grandes, de entre 31 y 50 casos, y entre el 59\,\% y el 86\,\% de las variables. Es decir, un bosque más
regularizado que el de partida, que generaliza mejor hacia años no vistos. La importancia fANOVA de Optuna lo
confirma: la proporción de variables explica el 75\,\% de la variación del resultado y el número de árboles solo
el 8\,\%, como se aprecia en la Figura~\ref{fig:optuna}.

\fig{fig12_optuna_historial_importancia}{1}{\caption{Historial de pruebas de Optuna, con las pruebas podadas, e
importancia de cada hiperparámetro.}\label{fig:optuna}}

\begin{lectura}
Los hiperparámetros son los ajustes internos del modelo, como el número de árboles o su profundidad. Optimizarlos es
similar a afinar un instrumento: el modelo mejoró de 0,882 a 0,891. Las cuatro estrategias llegaron prácticamente al
mismo resultado, pero Optuna lo hizo en menos tiempo porque deja de probar a mitad de camino las configuraciones que
no prometen. El mejor ajuste resultó ser un bosque prudente, con hojas grandes, que evita memorizar casos
particulares y por eso funciona mejor con años nuevos.
\end{lectura}

\subsection{Evaluación final en el periodo de prueba}

El modelo optimizado se reentrenó con los datos de 2003 a 2021 y se evaluó una sola vez con los de 2022 a 2026. El
umbral de decisión se eligió con las predicciones fuera de muestra de la validación cruzada, como el valor que
maximiza la F1, sin observar los datos de prueba. El modelo optimizado superó al de valores por defecto en todas las
métricas, según la Tabla~\ref{tab:test}.

\begin{table}[H]\centering\small
\caption{Desempeño en el periodo de prueba, 2022 a 2026, con 1\,878 emergencias}\label{tab:test}
\begin{tabular}{lcccccc}
\toprule
Modelo & Umbral & ROC AUC & PR AUC & F1 & Precisión & Sensibilidad \\
\midrule
Random Forest por defecto & 0,56 & 0,800 & 0,552 & 0,311 & 0,741 & 0,197 \\
Random Forest optimizado & 0,68 & \textbf{0,821} & \textbf{0,586} & \textbf{0,434} & 0,724 & \textbf{0,310} \\
Línea base al azar & & 0,500 & 0,216 & & 0,216 & \\
\bottomrule
\end{tabular}
\end{table}

\fig{fig13_evaluacion_prueba}{1}{\caption{Curvas ROC y de precisión y sensibilidad en el periodo 2022 a 2026, y
matriz de confusión del modelo optimizado.}\label{fig:test}}

El desempeño en prueba es menor que en validación, 0,82 frente a 0,89, y con el umbral elegido el modelo solo
detecta el 31\,\% de las emergencias graves, aunque acierta en el 72\,\% de las que señala, según la
Figura~\ref{fig:test}. La explicación se encuentra en los datos: entre 2022 y 2026 la proporción de emergencias
graves bajó al 21,6\,\%, frente al 31,2\,\% del entrenamiento, porque ahora se registran muchos más eventos
pequeños. El umbral, fijado con la proporción anterior, quedó alto. El modelo ordena bien las emergencias según su
riesgo, que es lo que mide el ROC AUC, pero el umbral debe recalibrarse con datos recientes antes de usarse. Por esa
razón, la parte prescriptiva se apoya en el orden de riesgo y no en un umbral fijo.

\begin{lectura}
En años que el modelo nunca vio, sigue distinguiendo bien los casos graves de los leves, en 82 de cada 100
comparaciones. Cuando señala una emergencia como grave, acierta siete de cada diez veces. Su punto débil es que deja
pasar muchas emergencias graves si se usa un corte fijo, porque el registro reciente incluye más eventos pequeños
que antes. La solución práctica no es descartar el modelo, sino usarlo para ordenar los reportes de mayor a menor
riesgo, como se propone en la analítica prescriptiva.
\end{lectura}

% =============================================================================
\section{Paso 5: \emph{reporting} con las cuatro analíticas}\label{sec:reporting}

\subsection{Analítica descriptiva: qué pasó}

La Figura~\ref{fig:tablero} reúne los indicadores principales y la Tabla~\ref{tab:kpi} los resume por periodos. En
23 años y medio se registraron 6\,593 emergencias, con 2,76 millones de personas impactadas por evento, 116 mil
damnificados, 177 fallecidos, 17\,382 viviendas destruidas sin contar 2026 y 778 mil camélidos perdidos. El periodo
de 2010 a 2015 fue el más dañino: con menos emergencias registradas reunió casi un millón de personas por evento, el
45\,\% de casos de alto impacto y 503 mil camélidos perdidos. El periodo de 2022 a 2026 tiene la mayor cantidad de
emergencias registradas, 1\,878, pero con menos personas impactadas por emergencia.

\begin{table}[H]\centering\small
\caption{Indicadores descriptivos por periodo}\label{tab:kpi}
\begin{tabular}{lrrrrrr}
\toprule
Periodo & Emergencias & Personas & Damnificados & Fallecidos & Camélidos & \% alto \\
 & & por evento & & & perdidos & impacto \\
\midrule
2003 a 2009 & 2\,013 & 613\,750 & 50\,772 & 57 & 13\,167 & 30 \\
2010 a 2015 & 1\,271 & 987\,615 & 50\,726 & 53 & 502\,765 & 45 \\
2016 a 2021 & 1\,431 & 819\,762 & 7\,373 & 31 & 99\,261 & 20 \\
2022 a 2026 & 1\,878 & 336\,415 & 7\,271 & 36 & 163\,087 & 22 \\
\midrule
Total & 6\,593 & 2\,757\,542 & 116\,142 & 177 & 778\,280 & 28 \\
\bottomrule
\end{tabular}
\end{table}

\fig{fig14_tablero_descriptivo}{1}{\caption{Tablero descriptivo de las emergencias en la región Puno entre 2003 y 2026.}\label{fig:tablero}}

\subsection{Analítica diagnóstica: por qué pasó}

Las pruebas de asociación con la condición de alto impacto, presentadas en la Tabla~\ref{tab:diag}, muestran que el
tipo de peligro es lo que más explica la gravedad de una emergencia, con una V de Cramér de 0,57 que indica una
asociación fuerte, seguido del mes, con 0,37. La provincia y la región natural, por sí solas, pesan mucho menos.
Las personas impactadas difieren de forma muy significativa entre peligros, según la prueba de Kruskal y Wallis,
con $H=904$ y $p<0{,}001$. En cambio, el número anual de emergencias no muestra una tendencia significativa entre
2003 y 2025, con $\rho=0{,}09$ de Spearman y $p=0{,}70$.

\begin{table}[H]\centering\small
\caption{Asociación entre cada variable y la condición de alto impacto mediante la prueba chi cuadrado}\label{tab:diag}
\begin{tabular}{lrrcc}
\toprule
Variable & $\chi^2$ & gl & Valor p & V de Cramér \\
\midrule
Tipo de peligro & 2\,140,2 & 9 & $<0{,}001$ & 0,57 \\
Mes & 907,2 & 11 & $<0{,}001$ & 0,37 \\
Provincia & 106,0 & 12 & $<0{,}001$ & 0,13 \\
Región natural & 16,8 & 1 & $<0{,}001$ & 0,05 \\
\bottomrule
\end{tabular}
\end{table}

Sin embargo, el mismo peligro no afecta por igual a todas las provincias, como se observa en la
Figura~\ref{fig:diag}. En Azángaro, el 80\,\% de las heladas registradas superó las 100 personas impactadas; en
Lampa, el 75\,\%, y en Chucuito, el 73\,\%. En Moho la proporción baja al 48\,\% y en San Román, una provincia más
urbana, al 54\,\%. La importancia por permutación del modelo final apunta en la misma dirección, según la
Figura~\ref{fig:perm}: la variable más importante, con amplia diferencia, es la tasa histórica de alto impacto de ese
mismo peligro en la provincia, cuya permutación reduce el ROC AUC en 0,15, seguida del tipo de peligro y del mes.

\fig{fig15_diagnostico_tasa_alto}{1}{\caption{Proporción de emergencias de alto impacto por provincia y peligro,
solo en las celdas con al menos 15 casos.}\label{fig:diag}}

\fig{fig16_importancia_permutacion}{0.78}{\caption{Importancia por permutación de las variables del modelo final,
medida en el periodo de prueba.}\label{fig:perm}}

\begin{lectura}
La gravedad de una emergencia depende sobre todo de qué peligro es y de en qué mes ocurre. Pero hay algo más: cada
provincia tiene su propia vulnerabilidad. Una helada en Azángaro termina siendo grave ocho de cada diez veces,
mientras que en Moho lo es la mitad de las veces. La mejor pista de lo que puede pasar es lo que ya pasó en ese mismo
lugar con ese mismo peligro, y eso es justamente lo que el modelo aprendió a usar.
\end{lectura}

\subsection{Analítica predictiva: qué puede pasar}

Además de clasificar cada emergencia, el modelo se utilizó para construir una matriz de riesgo por provincia y mes.
Para cada celda, el número esperado de emergencias de alto impacto en un año se calcula como

\begin{equation}
\widehat{A}_{p,m} = \overline{E}_{p,m} \times \overline{P}_{p,m},
\end{equation}

\noindent donde $\overline{E}_{p,m}$ es el promedio de emergencias de la provincia $p$ en el mes $m$ entre 2021 y
2025, y $\overline{P}_{p,m}$ es la probabilidad media de alto impacto que el modelo asigna a esas emergencias,
suavizada hacia el promedio de la provincia. En un año típico se esperan unas 76 emergencias graves en la región,
como muestra la Figura~\ref{fig:riesgo}. Melgar, con 12,5; Azángaro, con 11,1, y Puno, con 10,6, son las provincias
con más casos esperados, seguidas de Sandia, con 8,6; Moho y Yunguyo esperan alrededor de uno al año. El riesgo
cambia de lugar según la temporada: en lluvias encabezan Melgar, Azángaro y Sandia, y la celda más crítica del año es
Azángaro en marzo; en heladas encabezan Puno, Melgar y Azángaro, con picos en junio y julio.

\fig{fig17_matriz_riesgo}{1}{\caption{Matriz de riesgo: emergencias de alto impacto esperadas por provincia y mes en
un año típico.}\label{fig:riesgo}}

\begin{lectura}
Cada número de la matriz indica cuántas emergencias graves se esperan en esa provincia y ese mes en un año normal.
Las celdas más oscuras señalan dónde y cuándo debe estar lista la respuesta. Por ejemplo, Azángaro en marzo espera
2,6 emergencias graves, la cifra más alta del año, y Puno en junio y julio espera casi dos cada mes. Con esta matriz
se puede planificar con anticipación en lugar de reaccionar cuando la emergencia ya ocurrió.
\end{lectura}

\subsection{Analítica prescriptiva: qué conviene hacer}

La analítica prescriptiva recomienda acciones y no se limita a anticipar lo que puede ocurrir \citep{lepenioti2020}.
Se proponen dos decisiones.

\paragraph{Decisión A. Priorizar la evaluación de daños.} Los equipos que realizan la evaluación de daños son
escasos. Si cada día se ordenan los reportes según la probabilidad del modelo y se evalúan primero los de mayor
riesgo, la cobertura de las emergencias graves mejora de forma notable respecto de atender por orden de llegada,
como se observa en la Figura~\ref{fig:ganancia}. En el periodo de prueba, evaluar primero el 10\,\% de los reportes
con mayor riesgo cubre el 33\,\% de los casos graves; el 20\,\%, el 54\,\%; el 30\,\%, el 66\,\%, y la mitad, el
84\,\%. Por orden de llegada, esos porcentajes serían iguales a la capacidad. La regla no depende del umbral, por lo
que no la afecta el cambio en la proporción de casos graves.

\fig{fig18_curva_ganancia_priorizacion}{0.7}{\caption{Curva de ganancia en el periodo 2022 a 2026: proporción de
emergencias graves cubiertas según la capacidad de evaluación, priorizando con el modelo y por orden de
llegada.}\label{fig:ganancia}}

\begin{lectura}
Si el personal solo alcanza para evaluar primero uno de cada cinco reportes, y se eligen los que el modelo considera
más riesgosos, se atiende a más de la mitad de las emergencias graves. Atendiendo por orden de llegada se cubriría
solo una de cada cinco. Con el mismo personal, el modelo permite llegar antes a quienes más lo necesitan.
\end{lectura}

\paragraph{Decisión B. Pre posicionar kits de ayuda humanitaria.} Antes de cada temporada debe decidirse cuántos kits
de abrigo, techo y alimentos dejar en cada provincia. Se formuló un modelo de programación lineal resuelto con
PuLP. La demanda esperada de la provincia $p$ en la temporada, $d_p$, es el número de emergencias graves esperadas
según la matriz de riesgo, multiplicado por la mediana de personas impactadas en las emergencias graves de esa
provincia entre 2016 y 2025 y dividido entre cuatro personas por kit. El peso $w_p$, igual a uno más la proporción de
damnificados entre las personas impactadas de la provincia, otorga más prioridad a donde más personas pierden su
vivienda, y $c_p$ es la distancia aproximada por carretera desde un almacén central supuesto en Juliaca. Para
representar rendimientos decrecientes, la asignación de cada provincia se divide en tres tramos de un tercio de su
demanda, $x_p=y_{p1}+y_{p2}+y_{p3}$: el primer tercio cubre las necesidades más urgentes y vale el triple, el
segundo vale el doble y el tercero vale lo normal, es decir, $v=3,2,1$. Con un stock $S$, el modelo es

\begin{equation}
\min \sum_{p}\sum_{k=1}^{3} w_p\,v_k\left(\frac{d_p}{3}-y_{pk}\right) + 0{,}001\sum_p c_p\,x_p
\quad\text{sujeto a}\quad \sum_p x_p \le S,\qquad 0\le y_{pk}\le \frac{d_p}{3}.
\end{equation}

\noindent El primer término representa la necesidad no cubierta, ponderada; el segundo solo desempata a favor de
las provincias más cercanas. El modelo se resolvió para las dos temporadas y para tres niveles de stock, equivalentes
al 50, 75 y 100\,\% de la demanda. La demanda esperada asciende a 4\,887 kits en lluvias y 6\,302 en heladas, de
modo que las heladas exigen un 29\,\% más de abastecimiento.

\begin{table}[H]\centering\small
\caption{Kits asignados por provincia con un stock equivalente al 75\,\% de la demanda esperada}\label{tab:kits}
\begin{tabular}{lrrrrrr}
\toprule
 & \multicolumn{3}{c}{Lluvias, diciembre a marzo} & \multicolumn{3}{c}{Heladas, mayo a setiembre} \\
\cmidrule(lr){2-4}\cmidrule(lr){5-7}
Provincia & Demanda & Asignado & Cobertura & Demanda & Asignado & Cobertura \\
\midrule
Azángaro & 642 & 580 & 90\,\% & 600 & 600 & 100\,\% \\
Carabaya & 631 & 421 & 67\,\% & 1\,047 & 698 & 67\,\% \\
Chucuito & 298 & 199 & 67\,\% & 440 & 293 & 67\,\% \\
El Collao & 389 & 259 & 67\,\% & 337 & 225 & 67\,\% \\
Huancané & 319 & 213 & 67\,\% & 966 & 644 & 67\,\% \\
Lampa & 535 & 535 & 100\,\% & 631 & 631 & 100\,\% \\
Melgar & 914 & 609 & 67\,\% & 770 & 513 & 67\,\% \\
Moho & 22 & 15 & 67\,\% & 27 & 18 & 67\,\% \\
Puno & 521 & 347 & 67\,\% & 902 & 625 & 69\,\% \\
San Antonio de Putina & 92 & 61 & 67\,\% & 158 & 105 & 67\,\% \\
San Román & 230 & 230 & 100\,\% & 271 & 271 & 100\,\% \\
Sandia & 240 & 160 & 67\,\% & 147 & 98 & 67\,\% \\
Yunguyo & 54 & 36 & 67\,\% & 6 & 4 & 67\,\% \\
\midrule
Total & 4\,887 & 3\,665 & 75\,\% & 6\,302 & 4\,726 & 75\,\% \\
\bottomrule
\end{tabular}
\end{table}

Con un stock del 75\,\%, todas las provincias quedan cubiertas al menos en dos tercios y solo queda sin cubrir un
13\,\% de la necesidad ponderada, según la Tabla~\ref{tab:kits} y la Figura~\ref{fig:kits}. El último tercio se
completa primero en San Román, Lampa y Azángaro. Como los pesos de vulnerabilidad de las provincias son muy
parecidos, entre 1,02 y 1,12, en ese tramo decide la distancia, y esas son las provincias más cercanas al almacén de
Juliaca. Con un stock del 50\,\%, la necesidad no cubierta sube al 33\,\%. El precio sombra de la restricción de
stock es cercano a 2 cuando el stock es del 50\,\% y a 1 cuando es del 75\,\%, lo que significa que cada kit
adicional rinde el doble cuando el stock es escaso. La curva de sensibilidad de la Figura~\ref{fig:sens} muestra que
pasar del 50 al 75\,\% del stock reduce la necesidad no cubierta de 33 a 13\,\%, mientras que completar el 100\,\%
exige la misma cantidad de kits y rinde menos.

\fig{fig19_plan_kits}{1}{\caption{Asignación de kits y cobertura por provincia en las dos temporadas con un stock
equivalente al 75\,\% de la demanda esperada.}\label{fig:kits}}

\fig{fig20_sensibilidad_stock}{0.68}{\caption{Necesidad ponderada no cubierta según el stock disponible.}\label{fig:sens}}

\begin{lectura}
El modelo funciona como un reparto inteligente. Primero asegura a cada provincia dos tercios de lo que
probablemente necesitará y solo después completa a las que están más cerca del almacén. Con tres cuartas partes de
los kits se cubre el 87\,\% de la necesidad; conseguir el último cuarto cuesta lo mismo, pero aporta menos. El
precio sombra resume esta idea: cuando hay pocos kits, cada kit adicional vale el doble.
\end{lectura}

La recomendación que se desprende es pre posicionar al menos el 75\,\% de la demanda esperada antes de noviembre
para la temporada de lluvias, con prioridad en Melgar, Azángaro, Carabaya y Lampa, y antes de abril para la de
heladas, con prioridad en Carabaya, Huancané, Puno y Melgar.

% =============================================================================
\section{Tablero interactivo de inteligencia de negocios}\label{sec:tablero}

Para que los resultados puedan consultarse y explorarse sin leer el código, se desarrolló un tablero interactivo
en R con Shiny, publicado en \url{https://mlnp.shinyapps.io/emergencias-puno/}. El tablero conserva la identidad
visual del módulo elaborado en la semana 2 y organiza la información según las cuatro analíticas, de modo que el
usuario recorre las preguntas en el mismo orden que este informe. La Tabla~\ref{tab:pestanas} describe sus
pestañas.

\begin{table}[H]\centering\small
\caption{Pestañas del tablero interactivo y analítica a la que corresponden}\label{tab:pestanas}
\begin{tabularx}{\textwidth}{L{2.4cm}L{2.6cm}Y}
\toprule
Pestaña & Analítica & Contenido \\
\midrule
Resumen & Descriptiva & Indicadores clave, emergencias por año, frecuencia frente a impacto por peligro,
estacionalidad y distritos más afectados \\
Mapa & Descriptiva & Mapa distrital con cuatro métricas a elegir y detalle del distrito al hacer clic \\
Diagnóstico & Diagnóstica & Probabilidad de alto impacto por peligro y provincia, V de Cramér y calidad del registro \\
Modelo & Predictiva & Comparación de modelos, optimización de hiperparámetros, curva ROC, matriz de confusión e
importancia de variables \\
Predicción & Predictiva & Matriz de riesgo por provincia y mes, filtrable por temporada \\
Prescripción & Prescriptiva & Curva de priorización con capacidad ajustable y reparto de kits con stock ajustable \\
Datos & Todas & Tabla de emergencias filtradas con descarga en CSV \\
\bottomrule
\end{tabularx}
\end{table}

El tablero incorpora tres decisiones de diseño orientadas a que cualquier persona pueda interpretarlo. En primer
lugar, cada gráfico incluye debajo una explicación en lenguaje sencillo, que se recalcula cuando cambian los
filtros. En segundo lugar, los indicadores clave muestran, al pasar el cursor, qué significa la cifra y cómo debe
leerse. En tercer lugar, la pestaña de prescripción permite mover la capacidad de evaluación y el stock de kits para
observar cómo cambia la recomendación, lo que convierte el modelo de optimización en una herramienta de simulación
para el decisor. Para ello, la solución de PuLP se precalculó para niveles de stock entre el 30 y el 120\,\% de la
demanda, en pasos de cinco puntos. Las Figuras~\ref{fig:t1} a \ref{fig:t3} muestran tres de sus pestañas.

\fig{t_resumen}{1}{\caption{Pestaña Resumen del tablero, con los indicadores clave y la explicación sencilla bajo
cada gráfico.}\label{fig:t1}}

\fig{t_prediccion}{1}{\caption{Pestaña Predicción del tablero, con la matriz de riesgo por provincia y mes.}\label{fig:t2}}

\fig{t_prescripcion}{1}{\caption{Pestaña Prescripción del tablero, con los controles de capacidad y de stock.}\label{fig:t3}}

% =============================================================================
\section{Discusión y limitaciones}

El trabajo muestra que la calidad del registro condiciona cada nivel de la analítica. Sin revisar los datos se habría
concluido que 2019 fue un año tranquilo, que la ayuda humanitaria desapareció desde 2024 o que en el primer semestre
de 2026 se destruyeron más viviendas que en los 23 años anteriores. Las tres conclusiones habrían sido falsas, ya
que corresponden a problemas de registro. Este hallazgo coincide con lo que señalan \citet{sun2020} y
\citet{kyrkou2023}: en la gestión de emergencias, la principal limitación de la analítica no es el algoritmo, sino
los datos con que se alimenta.

En el modelado, la ganancia de la optimización fue modesta en validación, alrededor de 0,01 de ROC AUC, pero se
mantuvo en el periodo de prueba, donde alcanzó 0,02. Que las cuatro estrategias convergieran a valores casi iguales
indica que, con un espacio de cuatro hiperparámetros y presupuestos de 40 evaluaciones, la diferencia entre Optuna,
scikit-optimize e Hyperopt reside más en la eficiencia que en el resultado. Este comportamiento es consistente con
lo que describen \citet{bischl2023}: la ventaja de los métodos bayesianos y de la poda crece con la dimensión del
espacio y con el costo de cada evaluación.

En la parte prescriptiva, el aporte principal es conectar la predicción con la decisión. Tal como advierte la revisión de
\citet{sabbaghtorkan2020}, la mayoría de modelos de pre posicionamiento parten de escenarios de demanda supuestos;
aquí la demanda proviene de la matriz de riesgo estimada con datos, lo que vuelve la recomendación verificable y
actualizable cada año.

El trabajo tiene varias limitaciones. El umbral de 100 personas es una definición operativa elegida y no una norma
oficial. El modelo pierde desempeño en los años recientes porque cambió la proporción de casos graves, por lo que su
umbral de clasificación debe recalibrarse. Las personas impactadas son personas por evento y no personas distintas.
En el modelo prescriptivo, las distancias, el tamaño del kit y los pesos son supuestos razonables pero simplificados;
un uso real requeriría los costos de transporte, la capacidad de los almacenes y el inventario del INDECI. Por
último, el modelo no incorpora variables climáticas, como la temperatura mínima o la precipitación, que podrían
mejorar la anticipación.

% =============================================================================
\section{Conclusiones}

\begin{enumerate}
\item \textbf{Conjunto de datos.} El registro del SINPAD es una fuente local, abierta y rica para la inteligencia de
negocios en la gestión del riesgo de Puno, aunque presenta problemas de registro, como el subregistro de 2014 y
2019, la pérdida del registro de la ayuda y el error en las viviendas destruidas de 2026, que deben corregirse o
declararse antes de analizar.
\item \textbf{Análisis exploratorio.} Las heladas constituyen el problema central de la región: el 29\,\% de las
emergencias, pero el 77\,\% de las personas impactadas y el 76\,\% de los camélidos perdidos. El riesgo se concentra
en dos temporadas, de diciembre a marzo y de mayo a agosto, y en pocos distritos.
\item \textbf{Selección de modelos.} Con la información disponible al momento del reporte, un Random Forest anticipa si
una emergencia será de alto impacto con un ROC AUC de 0,882 en validación temporal.
\item \textbf{Optimización.} Optuna, scikit-optimize e Hyperopt mejoraron el modelo de manera similar, de 0,882 a 0,891 en
validación y de 0,800 a 0,821 en prueba. Optuna fue la estrategia más eficiente gracias a la poda y scikit-optimize
la que convergió en menos evaluaciones.
\item \textbf{\emph{Reporting}.} La analítica descriptiva identificó qué ocurrió y cuándo; la diagnóstica mostró que el
tipo de peligro y la vulnerabilidad histórica de cada provincia explican la gravedad; la predictiva estimó unas 76
emergencias graves por año, con Melgar, Azángaro y Puno a la cabeza, y la prescriptiva propuso dos decisiones
concretas. El tablero interactivo pone estos resultados al alcance de cualquier usuario.
\end{enumerate}

% =============================================================================
\section{Recomendaciones}

\begin{enumerate}
\item Fortalecer el llenado del SINPAD por parte de los gobiernos locales, en especial la evaluación de daños y el
registro de la ayuda entregada, ya que su ausencia impide medir la eficacia de la respuesta.
\item Incorporar el orden de riesgo del modelo en la programación diaria de los equipos de evaluación, empezando por las
provincias de Melgar, Azángaro y Puno.
\item Pre posicionar al menos el 75\,\% de los kits estimados antes de noviembre para la temporada de lluvias y antes de
abril para la de heladas.
\item Recalibrar el umbral del modelo cada año con las emergencias más recientes y actualizar la matriz de riesgo y el
plan de kits con el mismo cuaderno, que ya está programado para ello.
\item En trabajos futuros, integrar datos climáticos del SENAMHI y los costos logísticos reales del INDECI para mejorar
tanto la predicción como la optimización.
\end{enumerate}

% =============================================================================
\section*{Referencias}
% APA 7. Todas las referencias con DOI se verificaron en Crossref el 14/11/2026.
\renewcommand{\refname}{}
\setlength{\bibsep}{6pt}
\setlength{\bibhang}{1.27cm}
\vspace{-1.2cm}
\begin{thebibliography}{99}\small
\bibitem[Akiba et~al.(2019)]{akiba2019}
Akiba, T., Sano, S., Yanase, T., Ohta, T. y Koyama, M. (2019). Optuna: A next-generation hyperparameter
optimization framework. En \emph{Proceedings of the 25th ACM SIGKDD International Conference on Knowledge Discovery
\& Data Mining} (pp. 2623--2631). ACM. \url{https://doi.org/10.1145/3292500.3330701}

\bibitem[Balcik y Beamon(2008)]{balcik2008}
Balcik, B. y Beamon, B. M. (2008). Facility location in humanitarian relief. \emph{International Journal of
Logistics Research and Applications, 11}(2), 101--121. \url{https://doi.org/10.1080/13675560701561789}

\bibitem[Bergmeir y Benítez(2012)]{bergmeir2012}
Bergmeir, C. y Benítez, J. M. (2012). On the use of cross-validation for time series predictor evaluation.
\emph{Information Sciences, 191}, 192--213. \url{https://doi.org/10.1016/j.ins.2011.12.028}

\bibitem[Bergstra et~al.(2015)]{bergstra2015}
Bergstra, J., Komer, B., Eliasmith, C., Yamins, D. y Cox, D. D. (2015). Hyperopt: A Python library for model
selection and hyperparameter optimization. \emph{Computational Science \& Discovery, 8}(1), Artículo 014008.
\url{https://doi.org/10.1088/1749-4699/8/1/014008}

\bibitem[Bischl et~al.(2023)]{bischl2023}
Bischl, B., Binder, M., Lang, M., Pielok, T., Richter, J., Coors, S., Thomas, J., Ullmann, T., Becker, M.,
Boulesteix, A.-L., Deng, D. y Lindauer, M. (2023). Hyperparameter optimization: Foundations, algorithms, best
practices, and open challenges. \emph{WIREs Data Mining and Knowledge Discovery, 13}(2), Artículo e1484.
\url{https://doi.org/10.1002/widm.1484}

\bibitem[Breiman(2001)]{breiman2001}
Breiman, L. (2001). Random forests. \emph{Machine Learning, 45}(1), 5--32.
\url{https://doi.org/10.1023/A:1010933404324}

\bibitem[Cerqueira et~al.(2020)]{cerqueira2020}
Cerqueira, V., Torgo, L. y Mozetič, I. (2020). Evaluating time series forecasting models: An empirical study on
performance estimation methods. \emph{Machine Learning, 109}(11), 1997--2028.
\url{https://doi.org/10.1007/s10994-020-05910-7}

\bibitem[Chen y Guestrin(2016)]{chen2016}
Chen, T. y Guestrin, C. (2016). XGBoost: A scalable tree boosting system. En \emph{Proceedings of the 22nd ACM SIGKDD
International Conference on Knowledge Discovery and Data Mining} (pp. 785--794). ACM.
\url{https://doi.org/10.1145/2939672.2939785}

\bibitem[Delen y Ram(2018)]{delen2018}
Delen, D. y Ram, S. (2018). Research challenges and opportunities in business analytics. \emph{Journal of
Business Analytics, 1}(1), 2--12. \url{https://doi.org/10.1080/2573234X.2018.1507324}

\bibitem[Ghaffarian et~al.(2023)]{ghaffarian2023}
Ghaffarian, S., Taghikhah, F. R. y Maier, H. R. (2023). Explainable artificial intelligence in disaster risk
management: Achievements and prospective futures. \emph{International Journal of Disaster Risk Reduction, 98},
Artículo 104123. \url{https://doi.org/10.1016/j.ijdrr.2023.104123}

\bibitem[Kaufman et~al.(2012)]{kaufman2012}
Kaufman, S., Rosset, S., Perlich, C. y Stitelman, O. (2012). Leakage in data mining: Formulation, detection, and
avoidance. \emph{ACM Transactions on Knowledge Discovery from Data, 6}(4), 1--21.
\url{https://doi.org/10.1145/2382577.2382579}

\bibitem[Ke et~al.(2017)]{ke2017}
Ke, G., Meng, Q., Finley, T., Wang, T., Chen, W., Ma, W., Ye, Q. y Liu, T.-Y. (2017). LightGBM: A highly efficient
gradient boosting decision tree. En \emph{Advances in Neural Information Processing Systems, 30} (pp. 3146--3154).

\bibitem[Kyrkou et~al.(2023)]{kyrkou2023}
Kyrkou, C., Kolios, P., Theocharides, T. y Polycarpou, M. (2023). Machine learning for emergency management: A
survey and future outlook. \emph{Proceedings of the IEEE, 111}(1), 19--41.
\url{https://doi.org/10.1109/JPROC.2022.3223186}

\bibitem[Lepenioti et~al.(2020)]{lepenioti2020}
Lepenioti, K., Bousdekis, A., Apostolou, D. y Mentzas, G. (2020). Prescriptive analytics: Literature review and
research challenges. \emph{International Journal of Information Management, 50}, 57--70.
\url{https://doi.org/10.1016/j.ijinfomgt.2019.04.003}

\bibitem[Linardos et~al.(2022)]{linardos2022}
Linardos, V., Drakaki, M., Tzionas, P. y Karnavas, Y. L. (2022). Machine learning in disaster management: Recent
developments in methods and applications. \emph{Machine Learning and Knowledge Extraction, 4}(2), 446--473.
\url{https://doi.org/10.3390/make4020020}

\bibitem[Pedregosa et~al.(2011)]{pedregosa2011}
Pedregosa, F., Varoquaux, G., Gramfort, A., Michel, V., Thirion, B., Grisel, O., Blondel, M., Prettenhofer, P.,
Weiss, R., Dubourg, V., Vanderplas, J., Passos, A., Cournapeau, D., Brucher, M., Perrot, M. y Duchesnay, É.
(2011). Scikit-learn: Machine learning in Python. \emph{Journal of Machine Learning Research, 12}, 2825--2830.

\bibitem[Rawls y Turnquist(2010)]{rawls2010}
Rawls, C. G. y Turnquist, M. A. (2010). Pre-positioning of emergency supplies for disaster response.
\emph{Transportation Research Part B: Methodological, 44}(4), 521--534. \url{https://doi.org/10.1016/j.trb.2009.08.003}

\bibitem[Sabbaghtorkan et~al.(2020)]{sabbaghtorkan2020}
Sabbaghtorkan, M., Batta, R. y He, Q. (2020). Prepositioning of assets and supplies in disaster operations
management: Review and research gap identification. \emph{European Journal of Operational Research, 284}(1), 1--19.
\url{https://doi.org/10.1016/j.ejor.2019.06.029}

\bibitem[Saito y Rehmsmeier(2015)]{saito2015}
Saito, T. y Rehmsmeier, M. (2015). The precision-recall plot is more informative than the ROC plot when evaluating
binary classifiers on imbalanced datasets. \emph{PLOS ONE, 10}(3), Artículo e0118432.
\url{https://doi.org/10.1371/journal.pone.0118432}

\bibitem[Shahriari et~al.(2016)]{shahriari2016}
Shahriari, B., Swersky, K., Wang, Z., Adams, R. P. y de Freitas, N. (2016). Taking the human out of the loop: A
review of Bayesian optimization. \emph{Proceedings of the IEEE, 104}(1), 148--175.
\url{https://doi.org/10.1109/JPROC.2015.2494218}

\bibitem[Sun et~al.(2020)]{sun2020}
Sun, W., Bocchini, P. y Davison, B. D. (2020). Applications of artificial intelligence for disaster management.
\emph{Natural Hazards, 103}(3), 2631--2689. \url{https://doi.org/10.1007/s11069-020-04124-3}
\end{thebibliography}

\section*{Fuente de datos}
\renewcommand{\refname}{}
\vspace{-1.2cm}
\begin{thebibliography}{99}\small
\bibitem[INDECI(2026)]{indeci2026}
Instituto Nacional de Defensa Civil. (2026). \emph{Emergencias y daños a nivel nacional por departamento}
[Conjunto de datos: BD\_2003-2025\_EMERGENCIAS.csv, BD\_2026\_Semestre\_I.csv y diccionario de datos]. Plataforma
Nacional de Datos Abiertos. Licencia ODC BY. Última modificación: 30 de octubre de 2026. Consultado el 9 de
noviembre de 2026. \url{https://www.datosabiertos.gob.pe/dataset/emergencias-y-da\%C3\%B1os-nivel-nacional-por-departamento}
\end{thebibliography}

% =============================================================================
\section*{Anexo A. Glosario de términos}
\addcontentsline{toc}{section}{Anexo A. Glosario de términos}

\begin{table}[H]\centering\small
\begin{tabularx}{\textwidth}{L{3.6cm}Y}
\toprule
Término & Significado \\
\midrule
Emergencia de alto impacto & Emergencia con 100 o más personas impactadas. Es la condición que el modelo aprende a anticipar. \\
Personas impactadas & Suma de fallecidos, desaparecidos, lesionados, damnificados y afectados de una emergencia. \\
Validación cruzada temporal & Forma de evaluar un modelo en la que siempre se entrena con el pasado y se prueba con el
periodo siguiente, como ocurriría en la realidad. \\
Fuga de información & Error que consiste en usar para predecir un dato que todavía no se conoce en el momento de la
decisión, lo que hace que el modelo parezca mejor de lo que es. \\
ROC AUC & Probabilidad de que el modelo asigne más riesgo a una emergencia grave que a una leve tomadas al azar.
Vale 0,5 para el azar y 1 para un modelo perfecto. \\
PR AUC & Resume la relación entre la precisión y la sensibilidad. Es más exigente cuando los casos graves son minoría. \\
Precisión & De las emergencias que el modelo señala como graves, proporción que realmente lo fue. \\
Sensibilidad & De las emergencias realmente graves, proporción que el modelo logró detectar. \\
F1 & Promedio armónico de la precisión y la sensibilidad; equilibra ambos aciertos. \\
Puntaje de Brier & Error cuadrático medio de las probabilidades. Cuanto más bajo, mejor calibradas están. \\
Hiperparámetro & Ajuste del modelo que se fija antes de entrenarlo, como el número de árboles o su profundidad. \\
Poda & Técnica de Optuna que abandona una configuración a mitad de su evaluación cuando sus primeros resultados son
peores que la mediana de las anteriores. \\
V de Cramér & Medida de la fuerza de asociación entre dos variables categóricas, entre 0 y 1. \\
Importancia por permutación & Caída del desempeño del modelo cuando se desordenan los valores de una variable;
cuanto mayor es la caída, más importante es la variable. \\
Curva de ganancia & Proporción de casos graves cubiertos según la proporción de reportes atendidos en orden de riesgo. \\
Precio sombra & Mejora del objetivo de un modelo de optimización por cada unidad adicional de un recurso limitado,
en este caso, por cada kit adicional. \\
\bottomrule
\end{tabularx}
\end{table}

% =============================================================================
\section*{Anexo B. Archivos entregados y enlaces}

\begin{itemize}
\item \textbf{Tablero interactivo:} \url{https://mlnp.shinyapps.io/emergencias-puno/}
\item \textbf{Repositorio público:} \url{https://github.com/mlnp-dev/emergencias-puno-indeci}
\item \texttt{Semana4\_INDECI\_Puno.ipynb}: cuaderno con los cinco pasos, ejecutado, con el código, las salidas y la
interpretación de cada resultado, junto con su versión en HTML.
\item \texttt{datos/}: archivos originales del INDECI y la base limpia de Puno.
\item \texttt{figuras/}: las 20 figuras del análisis en formato PNG.
\item \texttt{resultados/}: tablas en CSV de calidad, selección de modelos, comparación de optimizadores, evaluación de
prueba, matriz de riesgo y plan de kits, además de los mejores hiperparámetros en JSON.
\item \texttt{shiny/}: código y datos del tablero interactivo; \texttt{capturas/}: imágenes de cada pestaña.
\end{itemize}

\end{document}
